% Part: first-order-logic
% Chapter: syntax-and-semantics
% Section: terms-formulas

\documentclass[../../../include/open-logic-section]{subfiles}

\begin{document}

\olfileid{fol}{syn}{frm}

\olsection{పదాలు మరియు \printtoken{P}{formula}}

ఒక మొదటిస్థాయి విధేయ భాష~$\Lang L$ను ఇచ్చిన తరువాత, $\Lang L$ యొక్క
ప్రాథమిక పదజాలం నుంచి నిర్మితమైన వ్యక్తీకరణలను నిర్వచించవచ్చు.
వీటిలో ముఖ్యంగా \emph{పదాలు}, \emph{\tetoken{సూత్రాలు}{!!{formula}s}} ఉంటాయి.

\begin{defn}[పదాలు]
\ollabel{defn:terms}
$\Lang L$ యొక్క \emph{పదాల} సమితి~$\Trm[L]$ను కింది విధంగా
ఆగమనాత్మకంగా నిర్వచిస్తాం:
\begin{enumerate}
\item ప్రతి \tetoken{చరం}{!!{variable}} ఒక పదం.
\item $\Lang L$లోని ప్రతి \tetoken{స్థిరసంకేతం}{!!{constant}} ఒక పదం.
\item $f$ ఒక $n$-స్థాన \tetoken{ప్రమేయ సంకేతం}{!!{function}}, $t_1$, \dots, $t_n$ పదాలు అయితే,
  $\Atom{f}{t_1, \ldots, t_n}$ ఒక పదం.
\tagitem{limitClause}{
  మరేదీ పదం కాదు.}{}
\end{enumerate}
ఏ \tetoken{చరాలనూ}{!!{variable}s} కలిగి లేని పదం \emph{సంవృత పదం}.
\end{defn}

\begin{explain}
మన భాషా నిర్దేశంలోనూ పదాల నిర్దేశంలోనూ \tetoken{స్థిరసంకేతాలను}{!!{constant}s} విడి సంకేత
వర్గంగా చేర్చాం; బదులుగా వాటిని శూన్యస్థాన \tetoken{ప్రమేయ సంకేతాలలో}{!!{function}s} చేర్చి
ఉండవచ్చు. అప్పుడు పదాల నిర్వచనంలో రెండవ నిబంధన అవసరం ఉండేది కాదు.
$n = 0$ అయితే $\Atom{f}{t_1, \ldots, t_n}$ను కేవలం $f$గానే
అర్థం చేసుకోవాలి.
\end{explain}

\begin{defn}[సూత్రాలు]
\ollabel{defn:formulas}
భాష~$\Lang L$ యొక్క \emph{\tetoken{సూత్రాలు}{!!{formula}s}} సమితి~$\Frm[L]$ను కింది
విధంగా ఆగమనాత్మకంగా నిర్వచిస్తాం:
\begin{enumerate}
\tagitem{prvFalse}{$\lfalse$ ఒక పరమాణు \tetoken{సూత్రం}{!!{formula}}.}{}

\tagitem{prvTrue}{$\ltrue$ ఒక పరమాణు \tetoken{సూత్రం}{!!{formula}}.}{}

\item $R$ అనేది $\Lang L$లోని $n$-స్థాన \tetoken{విధేయ సంకేతం}{!!{predicate}}, $t_1$, \dots,
  $t_n$ అనేవి $\Lang L$లోని పదాలు అయితే, $\Atom{R}{t_1,\ldots, t_n}$
  ఒక పరమాణు \tetoken{సూత్రం}{!!{formula}}.

\item $t_1$, $t_2$ అనేవి $\Lang L$లోని పదాలు అయితే,
  $\Atom{\eq}{t_1, t_2}$ ఒక పరమాణు \tetoken{సూత్రం}{!!{formula}}.

\tagitem{prvNot}{$!A$ \tetoken{ఒక సూత్రం}{!!a{formula}} అయితే, $\lnot !A$ కూడా
  \tetoken{ఒక సూత్రం}{!!a{formula}}.}{}

\tagitem{prvAnd}{$!A$, $!B$లు \tetoken{సూత్రాలు}{!!{formula}s} అయితే, $(!A \land
  !B)$ \tetoken{ఒక సూత్రం}{!!a{formula}}.}{}

\tagitem{prvOr}{$!A$, $!B$లు \tetoken{సూత్రాలు}{!!{formula}s} అయితే, $(!A \lor !B)$
  \tetoken{ఒక సూత్రం}{!!a{formula}}.}{}

\tagitem{prvIf}{$!A$, $!B$లు \tetoken{సూత్రాలు}{!!{formula}s} అయితే, $(!A \lif !B)$
  \tetoken{ఒక సూత్రం}{!!a{formula}}.}{}

\tagitem{prvIff}{$!A$, $!B$లు \tetoken{సూత్రాలు}{!!{formula}s} అయితే, $(!A \liff !B)$
  \tetoken{ఒక సూత్రం}{!!a{formula}}.}{}

\tagitem{prvAll}{$!A$ \tetoken{ఒక సూత్రం}{!!a{formula}}, $x$ \tetoken{ఒక చరం}{!!a{variable}} అయితే,
  $\lforall[x][!A]$ \tetoken{ఒక సూత్రం}{!!a{formula}}.}{}

\tagitem{prvEx}{$!A$ \tetoken{ఒక సూత్రం}{!!a{formula}}, $x$ \tetoken{ఒక చరం}{!!a{variable}} అయితే,
  $\lexists[x][!A]$ \tetoken{ఒక సూత్రం}{!!a{formula}}.}{}

\tagitem{limitClause}{మరేదీ \tetoken{ఒక సూత్రం}{!!a{formula}} కాదు.}{}
\end{enumerate}
\end{defn}

\begin{explain}
పదాల సమితి నిర్వచనం, \tetoken{సూత్రాలు}{!!{formula}s} సమితి నిర్వచనం \emph{ఆగమనాత్మక
నిర్వచనాలు}. సారాంశంగా, \tetoken{సూత్రాలు}{!!{formula}s} సమితిని అనంత సంఖ్యలో దశలుగా
నిర్మిస్తాం. ప్రారంభ దశలో పరమాణు సూత్రాలన్నీ సూత్రాలేనని ప్రకటిస్తాం;
ఇది నిర్వచనంలోని మొదటి కొన్ని సందర్భాలకు, అంటే
\iftag{prvTrue}{$\ltrue$, }{}%
\iftag{prvFalse}{$\lfalse$, }{}%
$\Atom{R}{t_1,\dots,t_n}$, $\Atom{\eq}{t_1,t_2}$ సందర్భాలకు సరిపోతుంది.
అందువల్ల ``పరమాణు \tetoken{సూత్రం}{!!{formula}}'' అంటే ఈ రూపాల్లో ఏదో ఒక రూపంలో ఉన్న
\tetoken{సూత్రం}{!!{formula}}.

నిర్వచనంలోని ఇతర సందర్భాలు, అప్పటికే నిర్మించిన \tetoken{సూత్రాలు}{!!{formula}s} నుంచి
కొత్త \tetoken{సూత్రాలను}{!!{formula}s} నిర్మించే నియమాలను ఇస్తాయి. రెండవ దశలో వాటిని
ఉపయోగించి పరమాణు \tetoken{సూత్రాలు}{!!{formula}s} నుంచి \tetoken{సూత్రాలను}{!!{formula}s} నిర్మించవచ్చు.
మూడవ దశలో పరమాణు సూత్రాల నుంచీ రెండవ దశలో లభించిన వాటి నుంచీ కొత్త
సూత్రాలను నిర్మిస్తాం; ఇలాగే కొనసాగిస్తాం. అటువంటి ఏదో ఒక దశలో
చివరకు నిర్మితమయ్యేదంతా \tetoken{సూత్రం}{!!{formula}}; మరేదీ కాదు.
\end{explain}

సంప్రదాయం ప్రకారం $\eq$ను దాని ఆర్గ్యుమెంట్ల మధ్య రాసి కుండలీకరణ
చిహ్నాలను విడిచిపెడతాం: $\eq[t_1][t_2]$ అనేది
$\Atom{\eq}{t_1,t_2}$కు సంక్షిప్త రూపం. అంతేకాక,
$\lnot \Atom{\eq}{t_1,t_2}$ను $\eq/[t_1][t_2]$గా సంక్షిప్తంగా
రాస్తాం. ద్విస్థానిక సంయోజకం~$\ast$ను ఉపయోగించి $!B$, $!C$ల నుంచి
నిర్మించిన $(!B \ast !C)$ సూత్రాన్ని రాసేటప్పుడు, తరచుగా వెలుపలి
కుండలీకరణ చిహ్నాల జతను విడిచిపెట్టి కేవలం~$!B \ast !C$ అని రాస్తాం.

\begin{intro}
కొన్ని తర్కశాస్త్ర గ్రంథాలు \tetoken{చరం}{!!{variable}}~$x$ తప్పనిసరిగా $!A$లో రావాలని
కోరుతాయి; అప్పుడు
\iftag{prvEx}{$\lexists[x][!A]$ }{}%
\iftag{notprvEx,notprvAll}{}{మరియు }%
\iftag{prvAll}{$\lforall[x][!A]$ }{}%
అనేది
\iftag{notprvEx,notprvAll}{\tetoken{ఒక సూత్రం}{!!a{formula}}}{\tetoken{సూత్రాలు}{!!{formula}s}}
గా లెక్కించాలి.
అలా కోరకపోయినా ఎలాంటి సమస్యా రాదు; పైగా విషయాలు సులభమవుతాయి.
\end{intro}

\begin{tagblock}{defNot,defOr,defAnd,defIf,defIff,defTrue,defFalse,defEx,defAll}
\begin{defn}
నిర్వచిత సంయోజకాలను ఉపయోగించి నిర్మించిన సూత్రాలను కింది విధంగా
అర్థం చేసుకోవాలి:

\begin{tagenumerate}{defTrue,defFalse,defNot,defOr,defAnd,defIf,defIff,defEx,defAll}

\tagitem{defTrue}{$\ltrue$ అనేది
  \iftag{prvFalse}{$\lnot\lfalse$కు}{$(!A \lor \lnot !A)$కు; ఇక్కడ
    $!A$ ఏదో ఒక స్థిరంగా ఎంచుకున్న పరమాణు \tetoken{సూత్రం}{!!{formula}}} సంక్షిప్త రూపం.}{}

\tagitem{defFalse}{$\lfalse$ అనేది
  \iftag{prvTrue}{$\lnot\ltrue$కు}{$(!A \land \lnot !A)$కు; ఇక్కడ
    $!A$ ఏదో ఒక స్థిరంగా ఎంచుకున్న పరమాణు \tetoken{సూత్రం}{!!{formula}}} సంక్షిప్త రూపం.}{}

\tagitem{defNot}{$\lnot !A$ అనేది $!A \lif \lfalse$కు సంక్షిప్త రూపం.}{}

\tagitem{defOr}{$!A \lor !B$ అనేది
  \iftag{prvAnd}{$\lnot(\lnot !A \land \lnot !B)$కు}{$\lnot !A \lif
    !B$కు} సంక్షిప్త రూపం.}{}

\tagitem{defAnd}{$!A \land !B$ అనేది
  \iftag{prvOr}{$\lnot(\lnot !A \lor \lnot !B)$కు}{$\lnot (!A \lif
    \lnot !B)$కు} సంక్షిప్త రూపం.}{}

\tagitem{defIf}{$!A \lif !B$ అనేది
  \iftag{prvOr}{$\lnot !A \lor !B$కు}{$\lnot (!A \land \lnot !B)$కు}
  సంక్షిప్త రూపం. \sourcecorrection{OLTEFOLSYN-001}{ఆంగ్ల మూలంలోని
  $\lnot !A \lor !B)$లో జతలేని చివరి కుండలీకరణ చిహ్నాన్ని తొలగించాం;
  నిర్వచిత సోపాధికానికి ఉద్దేశించిన $\lnot !A \lor !B$ రూపాన్ని
  యథాతథంగా ఉంచాం.}}{}

\tagitem{defIff}{$!A \liff !B$ అనేది $(!A \lif !B) \land (!B
  \lif !A)$కు సంక్షిప్త రూపం.}{}

\tagitem{defAll}{$\lforall[x][!A]$ అనేది $\lnot\lexists[x][\lnot !A]$కు సంక్షిప్త రూపం.}{}

\tagitem{defEx}{$\lexists[x][!A]$ అనేది $\lnot\lforall[x][\lnot !A]$కు సంక్షిప్త రూపం.}{}
\end{tagenumerate}
\end{defn}
\end{tagblock}

ఒక నిర్దిష్ట అనువర్తన భాషలో పనిచేసేటప్పుడు ద్విస్థానిక
\tetoken{విధేయ సంకేతాలు}{!!{predicate}s}, \tetoken{ప్రమేయ సంకేతాలను}{!!{function}s} వాటి సంబంధిత పదాల మధ్య తరచుగా రాస్తాం.
ఉదాహరణకు, అంకగణిత భాషలో $t_1 < t_2$, $(t_1 + t_2)$; సమితి సిద్ధాంత
భాషలో $t_1 \in t_2$. అంకగణిత భాషలో ఉత్తరవర్తి ప్రమేయాన్ని
సంప్రదాయంగా దాని ఆర్గ్యుమెంటు \emph{తరువాత} కూడా రాస్తాం:~$t'$.
అయితే అధికారికంగా ఇవి వరుసగా $\Atom{\Obj A^2_0}{t_1, t_2}$,
$\Obj f^2_0(t_1, t_2)$, $\Atom{\Obj A^2_0}{t_1, t_2}$,
$f^1_0(t)$లకు సంప్రదాయ సంక్షిప్త రూపాలు మాత్రమే.

\begin{defn}[వాక్యనిర్మాణ తాదాత్మ్యం]
$\ident$ సంకేతం సంకేతమాలల మధ్య వాక్యనిర్మాణ తాదాత్మ్యాన్ని
వ్యక్తపరుస్తుంది. అంటే, $!A$, $!B$లు సమాన పొడవు గల సంకేతమాలలు,
ప్రతి స్థానంలోనూ ఒకే సంకేతాన్ని కలిగి ఉన్నప్పుడు, అప్పుడు మాత్రమే
$!A \ident !B$.
\end{defn}

$\ident$ సంకేతానికి రెండు వైపులా సంయోజనం ద్వారా వచ్చిన సంకేతమాలలు
ఉండవచ్చు. ఉదాహరణకు, $!A \ident (!B \lor !C)$ అంటే: ఎడమ కుండలీకరణ
చిహ్నం, సంకేతమాల~$!B$, $\lor$ సంకేతం, సంకేతమాల~$!C$, కుడి కుండలీకరణ
చిహ్నాన్ని ఇదే క్రమంలో సంయోజించడం ద్వారా వచ్చిన సంకేతమాలే $!A$.
ఇది నిజమైతే, $!A$ మొదటి సంకేతం ఎడమ కుండలీకరణ చిహ్నమనీ, $!A$లో
$!B$ రెండవ సంకేతం వద్ద మొదలయ్యే ఉపసంకేతమాలగా ఉందనీ, ఆ ఉపసంకేతమాల
తరువాత~$\lor$ వస్తుందనీ, ఇలాగే మిగతావీ తెలుసుకుంటాం.

పదాలు, \tetoken{సూత్రాలను}{!!{formula}s} ప్రాథమిక అంశాల నుంచి ఆగమనాత్మక నిర్వచనాల ద్వారా
నిర్మిస్తాం కనుక, వాటి గురించి విషయాలను నిరూపించడానికి కింది ఆగమన
సూత్రాలను ఉపయోగించవచ్చు.

\begin{lem}[\emph{పదాలపై ఆగమన సూత్రం}]
    \ollabel{lem:trmind}
    $\Lang L$ ఒక మొదటిస్థాయి విధేయ భాష అనుకుందాం.
    ఏదో ధర్మం~$P$,
    %
    \begin{enumerate}
        \item ప్రతి \tetoken{చరం}{!!{variable}}~$v$కు వర్తిస్తే,
        %
        \item $\Lang L$లోని ప్రతి \tetoken{స్థిరసంకేతం}{!!{constant}}~$a$కు వర్తిస్తే, ఇంకా
        %
        \item $t_1$,~\dots, $t_n$లకు అది వర్తించినప్పుడల్లా,
        $f$ అనేది $\Lang L$లోని $n$-స్థాన \tetoken{ప్రమేయ సంకేతం}{!!{function}} అయితే,
        $f(t_1,\dotsc,t_n)$కు కూడా వర్తిస్తే
    \end{enumerate}
    ($t_1$,~\dots, $t_n$లు $\Lang{L}$లోని పదాలని అనుకుంటున్నాం),
    అప్పుడు $P$, $\Trm[L]$లోని ప్రతి పదానికీ వర్తిస్తుంది.
\end{lem}

\begin{prob}
    \olref[fol][syn][frm]{lem:trmind}ను నిరూపించండి.
\end{prob}

\begin{lem}[\emph{\tetoken{సూత్రాలపై}{!!{formula}s} ఆగమన సూత్రం}]
    \ollabel{thm:frmind}
    $\Lang L$ ఒక మొదటిస్థాయి విధేయ భాష అనుకుందాం.
    ఏదో ధర్మం~$P$ అన్ని పరమాణు \tetoken{సూత్రాలకు}{!!{formula}s} వర్తిస్తూ,
    కింది షరతులను కూడా తీరుస్తుందనుకుందాం:
    %
    \begin{enumerate}
        \tagitem{prvNot}{$!A$కు వర్తించినప్పుడల్లా
            $\lnot !A$కు వర్తిస్తుంది; }{}
        \tagitem{prvAnd}{$!A$,~$!B$లకు వర్తించినప్పుడల్లా
            $(!A \land !B)$కు వర్తిస్తుంది; }{}
        \tagitem{prvOr}{$!A$,~$!B$లకు వర్తించినప్పుడల్లా
            $(!A \lor !B)$కు వర్తిస్తుంది; }{}
        \tagitem{prvIf}{$!A$,~$!B$లకు వర్తించినప్పుడల్లా
            $(!A \lif !B)$కు వర్తిస్తుంది; }{}
        \tagitem{prvIff}{$!A$,~$!B$లకు వర్తించినప్పుడల్లా
            $(!A \liff !B)$కు వర్తిస్తుంది; }{}
        \tagitem{prvEx}{$!A$కు వర్తించినప్పుడల్లా
            $\lexists[x][!A]$కు వర్తిస్తుంది; }{}
        \tagitem{prvAll}{$!A$కు వర్తించినప్పుడల్లా
            $\lforall[x][!A]$కు వర్తిస్తుంది; }{}
  \end{enumerate}
  ($!A$, $!B$లు $\Lang{L}$లోని \tetoken{సూత్రాలు}{!!{formula}s} అని అనుకుంటున్నాం),
  అప్పుడు $P$, $\Frm[L]$లోని అన్ని సూత్రాలకూ వర్తిస్తుంది.
\end{lem}

\end{document}
